\section{Bilder und Tabellen und Codebeispiele}


In diesem Abschnitt werden Bilder und Tabellen eingebunden.
Auf Seite~\pageref{fig-example-picture} (unten) findet man in Abbildung
\ref{fig-example-picture} ein Beispiel.


Ebenfalls auf Seite~\pageref{tbl-cities} (aber hier oben)
gibt es eine Tabelle \ref{tbl-cities}.\\
\verb+\begin{figure}+ sowie \verb+\begin{table}+ erlauben für
eine automatische Nummerierung und Plazierung auf der Seite.
Ein Abbildungsverzeichnis wird automatisch erstellt.

\begin{b1important}{Zähler bei Floats}
\verb+\caption+ muss immer vor dem \verb+\label+
kommen
\end{b1important}


Und jetzt folgt ein Codebeispiel mit automatischer Nummerierung
und Auflistung im Index (Inhaltsverzeichnis).
\begin{lstlisting}[language={bash},
		caption={[Shell-Skript zur Demonstration]Shell-Skript das „Hallo Welt“ ausgibt},
		label={lst-helloworld},
		float={tbh},
		firstnumber={42},
		]
#!/bin/sh

echo "Hallo Welt!"

exit 0
\end{lstlisting}


Und ein generelles Codebeispiel; ohne im Index zu erscheinen:

\begin{Verbatim}[frame=leftline,numbers=left,label=server.conf,labelposition=topline]
[…]
client-connect <cmd>
client-disconnect <cmd>
[…]
\end{Verbatim}


%Place this floating object at the [b]ottom of the page, if possible
\begin{figure}[b]

	\centering
	\includegraphics[%draft,%
		width=5cm]{example}

	\caption{Ein Beispielbild}
	\label{fig-example-picture}
\end{figure}






% Place this floating object at the [t]op of the page; if possible
\begin{table}[t]
	\centering

	\captionabove{Städte, Länder und Flüsse}
	\label{tbl-cities}

	\begin{tabular}{l|c|r}
	\toprule
	Stadt	&	Land			&	Fluss \\
	\midrule
	München	&	Bayern			&	Isar \\
	Hamburg &	Niedersachsen	& Elbe \\
	Ingolstadt &	Bayern		&	Donau \\
	\bottomrule
	\end{tabular}

\end{table}




\begin{table}

	\centering
	\captionabove{Umfangreichere Tabelle}
	\label{tbl-extra}

	\begin{tabular}{@{}clc>{\columncolor{green3!30}}c>{\columncolor{green2!50}}c>{\columncolor{blue2!30}}c>{\columncolor{blue4!50}}c}
	\toprule

		& & & \multicolumn{2}{c}{\cellcolor{green1!80} Foobar}  & \multicolumn{2}{c}{\cellcolor{blue5}Acme}  \\
		\cmidrule(lr){4-7}
		\scriptsize Pos. & \scriptsize Kriterium & \scriptsize Gewichtung & \scriptsize Bewertung & \scriptsize Punkte & \scriptsize Bewertung  & \scriptsize Punkte  \\
		\midrule
		1 & Pretzel & 4 & 3         & 7          & 2  & 1  \\
		2 & Plunder & 3 & 5         & 4          & 5  & 6  \\
		3 & Sandwich & 8 & 1         & 1          & 7  & 3  \\
		4 & Fryd Butter Pastries & 2 & 9         & 1          & 6  & 2  \\
		5 & Bananas & 7 & 8         & 3          & 9  & 4  \\
		\midrule
		  & \multicolumn{1}{r}{Summe}       &  &  \cellcolor{white}        & 1337          & \cellcolor{white}    & 42  \\ \addlinespace
		  & \multicolumn{1}{r}{Platzierung} &  &  \cellcolor{white}        & \emph{1}          &  \cellcolor{white}   & \emph{2}  \\
	%1 & 2 & 3 & \cellcolor{red2} 4 & \cellcolor{blue5} 5 & \cellcolor{green1} 6 & 7 \\
	\bottomrule
	\end{tabular}%

\end{table}
% The @{} command kills the inter-column space and replaces it with whatever is between the curly braces.
% The >{} command executes the command in curly braces bevor the column element.  <{} is executed after
% The *{int}{spec} command repeats the columnspecification int times.
% The values of colors can have a {colorname!int} which indicates a lighter shade of the named color.
